\documentclass[11pt,letterpaper]{article}
\usepackage[T1]{fontenc}
\usepackage[utf8]{inputenc}
\usepackage[spanish,es-nodecimaldot,es-tabla,es-noshorthands]{babel}
\usepackage{lmodern}
\usepackage{microtype}
\usepackage[margin=2.2cm]{geometry}
\usepackage{amsmath,amssymb}
\usepackage{booktabs}
\usepackage{xcolor}
\usepackage{enumitem}
\usepackage{listings}
\usepackage{titlesec}
\usepackage[hidelinks]{hyperref}

\definecolor{marino}{RGB}{11,31,58}
\definecolor{verde}{RGB}{14,124,134}
\definecolor{coral}{RGB}{242,107,58}
\definecolor{fondo}{RGB}{238,243,248}
\setlist{itemsep=1pt, topsep=3pt}
\titleformat{\section}{\large\bfseries\color{marino}}{\thesection.}{0.6em}{}
\lstset{language=Python, basicstyle=\ttfamily\small, backgroundcolor=\color{fondo}, frame=none,
  columns=fullflexible, keepspaces=true, xleftmargin=1em, framexleftmargin=0.6em}

\newcommand{\vu}{\mathbf{u}}
\newcommand{\Rey}{\mathit{Re}}
% --- diálogo
\newcommand{\habla}[3]{\par\noindent\hangindent=4.6em\hangafter=1%
  \makebox[4.6em][l]{\textbf{\textcolor{#1}{#2:}}}#3\par\medskip}
\newcommand{\luis}[1]{\habla{marino}{Luis}{#1}}
\newcommand{\marco}[1]{\habla{verde}{Marco}{#1}}
\newcommand{\accion}[1]{\par\noindent{\small\itshape\color{gray}[#1]}\par\medskip}
\newcommand{\lee}[1]{\par\noindent{\small\color{coral}\textbf{Se lee:} #1}\par\medskip}
\newcommand{\diapo}[4]{\vspace{0.6em}\par\noindent\colorbox{fondo}{\parbox{\dimexpr\linewidth-2\fboxsep}{%
  \textbf{\color{marino}Diapositiva #1 · #2}\hfill{\small #3 · #4}}}\par\medskip}

\title{\textbf{Guion de la exposición}\\[2pt]\large Flujo laminar sobre un escalón descendente}
\author{Luis Antonio Pérez Hernández \and Marco Antonio Navarrete Baños\\[4pt]
  \small Mecánica de Fluidos · Proyecto final · Prof.\ Anuar Sánchez}
\date{}

\begin{document}
\maketitle
\vspace{-1.5em}

\section*{Cómo usar este guion}
\begin{itemize}
  \item El diálogo sale del reporte (\texttt{reporte/reporte.pdf}); cada diapositiva indica la sección del reporte de
    donde viene, por si quieren repasar más a fondo.
  \item \textbf{No es para leerlo.} La rúbrica penaliza la ``lectura literal de diapositivas''. Úsenlo para ensayar:
    lean una vez, luego digan cada bloque con sus palabras mirando sólo la diapositiva. Lo esencial está en
    \textbf{negritas}.
  \item Duración planeada: \textbf{18 min} (la rúbrica pide de 15 a 20). Si van tarde, recorten las frases marcadas
    como \emph{opcional}.
  \item Cada quien expone una parte técnica completa (requisito de la participación individual):
    \textbf{Luis}, la formulación física y la verificación; \textbf{Marco}, el método numérico, el código y la
    demostración.
  \item Las notas del orador del PowerPoint tienen una versión corta de este guion (vista del moderador).
\end{itemize}

\begin{table}[h]
\centering\small
\caption{Reparto y tiempos.}
\begin{tabular}{@{}rlll@{}}
\toprule
& Diapositiva & Quién expone & Tiempo\\
\midrule
1 & Portada & Ambos & 0:30\\
2 & El fenómeno y el caso base & Luis & 1:15\\
3 & Ecuaciones de gobierno & Luis & 1:45\\
4 & Condiciones iniciales y de frontera & Luis & 1:30\\
5 & Método de proyección & Marco & 1:45\\
6 & Discretización espacial & Marco & 1:45\\
7 & Código I: Poisson y fronteras & Marco & 1:15\\
8 & Código II: ciclo de tres pasos & Marco & 1:00\\
-- & \textbf{Demostración en vivo (Jupyter)} & Marco opera, Luis comenta & 2:30\\
9 & Estabilidad y convergencia & Luis & 1:30\\
10 & Análisis paramétrico: viscosidad & Marco & 1:30\\
11 & Validación & Luis & 1:00\\
12 & Conclusiones & Ambos & 0:45\\
\midrule
& & \textbf{Total} & \textbf{18:00}\\
\bottomrule
\end{tabular}
\end{table}

\section*{Antes de empezar (5 minutos antes)}
\begin{enumerate}
  \item Abrir \texttt{presentacion/presentacion.pptx} en modo presentación con vista del moderador (notas visibles).
  \item Abrir \texttt{escalon\_descendente.ipynb} en Jupyter, poner \texttt{HACER\_VIDEOS = False} en el bloque~0 y
    ejecutar las celdas \textbf{desde el inicio hasta la sección 4 (caso base) inclusive}: tarda menos de 10~s. Así las
    funciones ya están cargadas para la demostración.
  \item Dejar Jupyter en la celda de parámetros, en otra ventana, para cambiar rápido.
\end{enumerate}

% =====================================================================
\section{Presentación}
\diapo{1}{Portada}{Ambos}{0:30}

\luis{Buenos días. Somos Luis Antonio Pérez y Marco Antonio Navarrete. Nuestro proyecto es la simulación del
\textbf{flujo laminar sobre un escalón descendente}: resolvimos las ecuaciones de Navier–Stokes en dos dimensiones con un
código propio en Python.}

\marco{La exposición sigue el orden en que se construye una simulación: \textbf{la física y las ecuaciones, el método
numérico, el código funcionando y, al final, cómo sabemos que el resultado es correcto}. Luis empieza con el fenómeno.}

% =====================================================================
\section{Problema físico y formulación (Luis)}
\diapo{2}{El fenómeno y el caso base}{Luis · reporte §1 y §4}{1:15}

\luis{Un fluido viaja por un canal y encuentra un escalón: el canal se \textbf{ensancha al doble}. Al ensancharse, el
fluido se frena y, como en Bernoulli, \textbf{la presión sube}: es un gradiente de presión adverso,
$\partial p/\partial x>0$. El fluido pegado a la pared es lento y no tiene inercia para avanzar contra ese aumento de
presión, así que \textbf{se separa en la esquina}.}
\accion{Señalar la figura de arriba (velocidad).}
\luis{Debajo queda un \textbf{remolino} donde la velocidad es negativa, y más adelante el flujo vuelve a pegarse al piso
en el \textbf{punto de reacoplamiento}, donde el esfuerzo cortante en la pared es cero. En la figura de abajo se ve la
presión subiendo después del escalón. En nuestro caso base, con Reynolds 200, el remolino mide \textbf{4.91 alturas de
escalón}.}

\marco{¿Y por qué es un caso de prueba tan usado?}

\luis{Porque la geometría es simple pero tiene todos los mecanismos de un flujo separado: separación, capa de corte,
recirculación y reacoplamiento. Y hay experimentos, como los de Armaly en 1983, y muchas simulaciones para comparar.}

% ---------------------------------------------------------------------
\diapo{3}{Ecuaciones de gobierno}{Luis · reporte §2.1–2.3}{1:45}

\luis{Las ecuaciones son las de clase. La primera es \textbf{continuidad}: como el fluido es incompresible, la
divergencia de la velocidad es cero; no se crea ni se pierde masa.}
\lee{``divergencia de $\vu$ igual a cero''.}
\luis{Las otras dos son \textbf{Navier–Stokes} en $x$ y en $y$: aceleración local más convección igual a menos el
gradiente de presión más la difusión viscosa. Las escribimos \textbf{adimensionales}: longitudes entre $H$, velocidades
entre la velocidad media de entrada $U$ y presión entre $\rho U^2$. Así queda \textbf{un solo parámetro, el número de
Reynolds}, $\Rey=UH/\nu$. Con $U$ y $H$ iguales a uno, $\nu=1/\Rey$: \textbf{variar la viscosidad es lo mismo que variar
Reynolds}.}
\lee{``$\partial u/\partial t$ más $\partial(u^2)/\partial x$ más $\partial(uv)/\partial y$ igual a menos
$\partial p/\partial x$ más uno sobre Re por el laplaciano de $u$''.}
\luis{Usamos la \textbf{forma conservativa}: en lugar de $u\,\partial u/\partial x$ escribimos $\partial(u^2)/\partial x$.
Es lo mismo gracias a continuidad, porque $(\vu\cdot\nabla)\vu=\nabla\cdot(\vu\vu)-\vu(\nabla\cdot\vu)$ y el último
término es cero. La ventaja es que al discretizar, \textbf{el momento que sale de una celda es exactamente el que entra a la
vecina}.}
\luis{Los supuestos: flujo 2D, incompresible, newtoniano, isotérmico y \textbf{laminar}, porque trabajamos hasta Reynolds
600 y la transición en este flujo es alrededor de 1200. La gravedad la omitimos porque sólo agrega una presión
hidrostática que no cambia el movimiento.}

\marco{Entonces la presión no tiene una ecuación propia.}

\luis{Exacto: en incompresible no hay ecuación de estado. \textbf{La presión es la que haga falta para que se cumpla
continuidad}. Eso es justo lo que resuelve el método que va a explicar Marco.}

% ---------------------------------------------------------------------
\diapo{4}{Condiciones iniciales y de frontera}{Luis · reporte §2.4}{1:30}

\luis{La \textbf{condición inicial} es el fluido en reposo; avanzamos en el tiempo hasta que el flujo deja de cambiar, que
es el estado estacionario.}
\luis{En la \textbf{entrada}, a dos alturas antes del escalón, ponemos el \textbf{perfil parabólico de Poiseuille}, que
es la solución exacta para un canal largo: $u=6U\eta(1-\eta)$. El 6 sale de pedir que la velocidad media sea $U$, porque
la integral de $6\eta(1-\eta)$ de 0 a 1 vale uno; la máxima es $1.5U$. \emph{Opcional:} el canal de entrada deja que el
perfil se deforme antes de la esquina; sin él, la burbuja sale cerca de 8\,\% más larga.}
\luis{En \textbf{todas las paredes sólidas}, piso, techo y las dos caras del escalón, \textbf{no deslizamiento}:
$u=v=0$. En la \textbf{salida} suponemos flujo desarrollado, derivadas en $x$ iguales a cero, y fijamos $p=0$ como
referencia. Por eso el canal debe ser más largo que el remolino: usamos hasta 25 alturas.}
\luis{Para la presión, en las paredes la condición es de \textbf{Neumann}, $\partial p/\partial n=0$, que sale de evaluar
Navier–Stokes donde la velocidad es cero. Y el $p=0$ de la salida hace que el sistema de la presión tenga solución única.
Con esto el problema está planteado matemáticamente. Marco, ¿cómo lo resolvemos?}

% =====================================================================
\section{Método numérico y código (Marco)}
\diapo{5}{Método de proyección de Chorin}{Marco · reporte §3.2}{1:45}

\marco{El problema es que tenemos tres incógnitas, $u$, $v$ y $p$, y la presión no tiene ecuación de evolución. El
\textbf{método de proyección de Chorin}, de 1968, separa cada paso de tiempo en tres pasos.}
\marco{\textbf{Paso uno:} calculamos una \textbf{velocidad de prueba} $\vu^*$ con convección y difusión, pero sin la
presión. Para avanzar en el tiempo usamos \textbf{Adams–Bashforth}: la aceleración a la mitad del paso se estima con la
actual y la anterior, tres medios de la de ahora menos un medio de la anterior.}
\lee{``$\vu^*$ igual a $\vu^n$ más $\Delta t$ por tres medios de $A^n$ menos un medio de $A^{n-1}$''.}
\marco{\textbf{Paso dos:} esa velocidad no cumple continuidad. Queremos que $\vu^{n+1}=\vu^*-\Delta t\nabla p$ tenga
divergencia cero. Aplicamos divergencia a esa ecuación y queda la \textbf{ecuación de Poisson para la presión}: el
laplaciano de $p$ igual a la divergencia de $\vu^*$ entre $\Delta t$.}
\lee{``laplaciano de $p^{n+1}$ igual a divergencia de $\vu^*$ entre $\Delta t$''.}
\marco{\textbf{Paso tres:} corregimos la velocidad con el gradiente de esa presión, y ahora sí la divergencia es cero.
Si sumas el paso uno y el tres, \textbf{recuperas Navier–Stokes discretizada en el tiempo}.}

\luis{¿Por qué Adams–Bashforth y no Euler, que es más simple?}

\marco{Lo probamos: con Euler y convección centrada los errores crecen un poco en cada paso y \textbf{la solución explota
a los 706 pasos}, aunque el paso de tiempo sea la mitad del límite. Adams–Bashforth cuesta lo mismo, porque la
aceleración anterior ya está calculada, y es estable. Lo van a ver en la demostración.}

% ---------------------------------------------------------------------
\diapo{6}{Discretización espacial}{Marco · reporte §3.1 y §3.3}{1:45}

\marco{Para resolverlo en la computadora partimos el canal en \textbf{celdas cuadradas} y cambiamos cada derivada por una
diferencia. Usamos una \textbf{malla escalonada}: la presión vive en el centro de cada celda y las velocidades en las
caras.}
\accion{Señalar el dibujo de la malla.}
\marco{¿Por qué? Porque así \textbf{la velocidad de cada cara queda justo entre dos presiones}, y la fuerza de presión
que la empuja es simplemente la diferencia de esas dos. Si todo estuviera en el mismo punto, la presión podría oscilar
en tablero de ajedrez sin que el método lo notara.}
\marco{La \textbf{difusión} es el laplaciano de cinco puntos: los cuatro vecinos menos cuatro veces el centro, entre
$\Delta x^2$. La \textbf{convección} va en forma conservativa, lo que sale menos lo que entra, y el valor en cada cara es
el \textbf{promedio de los dos vecinos}. Las dos son de \textbf{segundo orden}: al duplicar la malla el error baja unas
cuatro veces. \emph{Opcional:} un esquema \emph{upwind} sería más estable, pero agrega una viscosidad artificial que
acorta la burbuja.}
\marco{La ecuación de Poisson queda así: para cada celda, la suma de las diferencias con sus vecinos es igual al lado
derecho. Si el vecino es pared, no se suma, que es la condición de Neumann. En la salida la cara está a media celda con
$p=0$, y eso da un coeficiente de menos dos. El escalón es simplemente un \textbf{bloque de celdas sólidas}.}

% ---------------------------------------------------------------------
\diapo{7}{Código I: matriz de Poisson y fronteras}{Marco · reporte anexo B}{1:15}

\marco{Esto es el código real del notebook. Cada celda se numera como $k=i\cdot N_y+j$, y \textbf{cada fila de la matriz
es la ecuación de una celda}. Si la celda es sólida, la ecuación es trivial. Para cada uno de los cuatro vecinos: si
es la salida, menos dos en la diagonal; si es fluido, más uno al vecino y menos uno a la diagonal. Como \textbf{la matriz
no cambia en el tiempo, se factoriza una sola vez} con LU y en cada paso sólo se resuelve.}
\marco{Abajo están las fronteras. Usamos \textbf{celdas fantasma}: una fila extra fuera de la pared. Si el fantasma vale
menos el valor del primer nodo, \textbf{el promedio en la pared da cero}.}

\luis{O sea que el no deslizamiento no se pone en un nodo, sino a la mitad entre el fantasma y el primer nodo real.}

\marco{Exacto.}

% ---------------------------------------------------------------------
\diapo{8}{Código II: el ciclo de tres pasos}{Marco · reporte §3.2 y §3.4}{1:00}

\marco{Y este es el corazón del programa. Cada línea es una ecuación de las que vimos: calculamos la aceleración, el
\textbf{paso uno} con Adams–Bashforth, la divergencia de $\vu^*$, el \textbf{paso dos} resolviendo Poisson y el
\textbf{paso tres} restando las diferencias de presión en cada cara. Al final medimos cuánto cambió la velocidad: cuando el
cambio por unidad de tiempo es menor que $10^{-4}$, \textbf{llegamos al estado estacionario}. Ahora lo vemos correr.}

% =====================================================================
\section{Demostración en vivo (Marco opera, Luis comenta)}
\diapo{--}{Jupyter}{Ambos · 50\,\% de la exposición grupal}{2:30}

\accion{Cambiar a la ventana de Jupyter. Celdas del bloque 0 a 4 ya ejecutadas.}

\marco{Los parámetros están al inicio. Vamos a bajar la viscosidad: cambio Reynolds de 200 a \textbf{400}.}
\accion{Escribir \texttt{RE = 400}; ejecutar la celda de parámetros y la del caso base (Ctrl+Enter en cada una).}
\luis{Mientras corre: está haciendo unos \textbf{8\,300 pasos de tiempo}, y en cada uno los tres pasos del método. Tarda
unos 8 segundos.}
\accion{Salida esperada: \texttt{converge en t = 138.2 (8294 pasos) | x\_r/S = 8.26}; divergencia del orden de
$10^{-15}$.}
\marco{Convergió: el remolino pasó de 4.91 a \textbf{8.26 alturas de escalón}, y la divergencia sigue siendo
prácticamente cero.}
\accion{Ejecutar la celda de figuras.}
\luis{Se ve el remolino mucho más largo, y aparece una \textbf{pequeña burbuja en el techo}, entre 4.2 y 4.5 alturas: con
menos viscosidad también se separa la capa límite de arriba.}

\marco{Ahora probamos la estabilidad. Primero el paso de tiempo justo en el límite de Courant:}
\begin{lstlisting}
simular(200, dt=(H/20)/1.5, t_max=150)
\end{lstlisting}
\accion{Salida esperada: \texttt{diverge en t = 1.2 (36 pasos)}.}
\luis{\textbf{Explotó en 36 pasos}: el fluido avanza más de una celda por paso y los errores crecen.}
\marco{Y ahora con el paso normal, pero con Euler en lugar de Adams–Bashforth:}
\begin{lstlisting}
simular(200, ab2=False, t_max=150)
\end{lstlisting}
\accion{Salida esperada: \texttt{diverge en t = 11.8 (706 pasos)}.}
\marco{\textbf{Explota a los 706 pasos}, aunque el paso de tiempo sea la mitad del límite. Por eso usamos
Adams–Bashforth.}
\accion{Regresar \texttt{RE = 200} y volver a la presentación.}

% =====================================================================
\section{Verificación y resultados}
\diapo{9}{Estabilidad y convergencia}{Luis · reporte §3.4, §5.2 y §5.3}{1:30}

\luis{Lo que acabamos de ver es la \textbf{condición CFL}: en un paso de tiempo el fluido no debe avanzar más de una
celda, $\Delta t\le\Delta x/|u|_{\max}$. Con celdas de 0.05 y velocidad máxima 1.5, el límite es 0.033. El límite por
difusión es 0.125, así que manda el CFL; nosotros usamos la mitad.}
\lee{``delta t menor o igual a delta x entre la velocidad máxima''.}
\luis{En la tabla: con 0.25, 0.5 y 0.8 veces el límite, el resultado es \textbf{idéntico}, 4.914; sólo cambia el número de
pasos. En el límite explota.}
\luis{Para la \textbf{malla} corrimos 10, 20 y 40 celdas por altura: 4.754, 4.914 y 4.946. El cambio se reduce unas cinco
veces cada vez que duplicamos la malla. Con eso calculamos el \textbf{orden observado}: logaritmo base dos del cociente
de los cambios, \textbf{2.3}, que coincide con el segundo orden del esquema. Con la extrapolación de Richardson el valor
con malla infinita es 4.95, así que \textbf{con 20 celdas el error es de 0.8\,\%}.}

\marco{¿Y la masa sí se conserva?}

\luis{Sí: la divergencia final es del orden de $10^{-15}$, que es el cero de la computadora.}

% ---------------------------------------------------------------------
\diapo{10}{Análisis paramétrico: viscosidad}{Marco · reporte §4.1 y §5.1}{1:30}

\accion{Reproducir el video.}
\marco{Aquí están tres arranques desde el reposo, con \textbf{Reynolds 100, 300 y 600}. La tinta son partículas que el
flujo arrastra. Con Reynolds 100 todo se calma pronto; con 600 se forman varios vórtices que viajan y se juntan, y aparece
la burbuja del techo.}
\marco{En la tabla: el remolino pasa de \textbf{2.84 a 10.06 alturas de escalón}. ¿Por qué crece? Con menos viscosidad,
la difusión tarda más en transmitir el momento del chorro hacia el piso, mientras la convección lo lleva más lejos, así
que se vuelve a pegar más adelante. Crece casi lineal hasta Reynolds 400; después más lento, porque \textbf{la burbuja del
techo estrecha el paso y empuja el chorro hacia abajo}.}

\luis{Y el tiempo para llegar al estacionario también crece, de 20 a 289 unidades: el transitorio lo controla la
difusión viscosa.}

% ---------------------------------------------------------------------
\diapo{11}{Validación}{Luis · reporte §6}{1:00}

\luis{Comparamos con simulaciones 2D publicadas del mismo problema. Entre Reynolds 200 y 600 la diferencia es de
\textbf{2\,\% o menos}; en Reynolds 100 es 5\,\%. Como nuestro error de malla es menor a 1\,\%, esas diferencias son del
tamaño de las diferencias entre los propios estudios publicados, que usan canales de entrada y dominios distintos. Y
reproducimos el remolino del techo a partir de Reynolds 400, como en la literatura.}

% ---------------------------------------------------------------------
\diapo{12}{Conclusiones}{Ambos}{0:45}

\luis{En resumen: planteamos continuidad y Navier–Stokes con sus condiciones de frontera, y verificamos el resultado:
\textbf{el método es estable dentro del CFL y de segundo orden, con error de malla menor a 1\,\%}.}

\marco{Lo resolvimos con el método de proyección en malla escalonada, \textbf{la masa se conserva a precisión de
máquina}, reproducimos la separación, la recirculación y el reacoplamiento, y coincidimos con la literatura dentro de
2\,\%. Muchas gracias; ¿tienen preguntas?}

% =====================================================================
\newpage
\section*{Preguntas individuales probables}
La rúbrica da el 40\,\% de la calificación individual a una pregunta del profesor, y no se puede pedir ayuda al compañero.
Cada quien debe dominar \textbf{todo}, pero sobre todo su parte. Más preguntas en
\texttt{guia/guia\_explicacion.pdf}, sección C.

\subsection*{Luis (formulación, estabilidad, validación)}
\begin{description}[leftmargin=1em, style=nextline]
  \item[¿Por qué se separa el flujo?] El canal se ensancha, el fluido se frena y la presión sube: gradiente adverso. El
    fluido junto a la pared es lento y no puede avanzar contra ese aumento de presión.
  \item[¿Por qué la velocidad máxima final es 0.75?] Por continuidad, el mismo caudal pasa por un canal el doble de alto:
    la media baja de 1 a 0.5 y el máximo de la parábola es $1.5\times0.5=0.75$.
  \item[¿Qué pasa si cambias la condición de salida?] La salida supone flujo desarrollado. Si el canal fuera más corto que
    el remolino, la condición estaría mal y el resultado cambiaría; por eso usamos hasta 25 alturas.
  \item[¿Cómo calculaste el orden de convergencia?] Con tres mallas de razón 2:
    $p=\log_2\big[(x_{20}-x_{10})/(x_{40}-x_{20})\big]=\log_2(0.160/0.032)\approx2.3$.
  \item[¿Qué es el número de Reynolds aquí?] $\Rey=UH/\nu$, con la velocidad media de entrada y la altura del canal
    ensanchado, igual al diámetro hidráulico del canal de entrada, $2h$.
\end{description}

\subsection*{Marco (método, código, demostración)}
\begin{description}[leftmargin=1em, style=nextline]
  \item[¿De dónde sale la ecuación de Poisson?] De aplicar divergencia a $\vu^{n+1}=\vu^*-\Delta t\nabla p$ y pedir
    $\nabla\cdot\vu^{n+1}=0$.
  \item[¿Por qué la matriz se factoriza una sola vez?] Porque sólo depende de la malla y de las condiciones de frontera, no
    del tiempo ni de la velocidad; en cada paso sólo cambia el lado derecho.
  \item[¿Qué es una celda fantasma?] Una fila fuera del dominio cuyo valor es el opuesto del primer nodo, para que el
    promedio sobre la pared dé cero (no deslizamiento).
  \item[¿Por qué malla escalonada?] Para que cada velocidad quede entre dos presiones y la fuerza de presión use vecinos
    directos; así no aparecen oscilaciones en tablero de ajedrez.
  \item[¿Cómo sabe el código que llegó al estado estacionario?] Cuando $\max|\vu^{n+1}-\vu^n|/\Delta t<10^{-4}$; con
    $10^{-5}$ el resultado cambia menos de 0.3\,\%.
\end{description}

\end{document}
